% ==============================================================================
% Introduction
% ==============================================================================

\section{Introduction}
\label{sec:introduction}

Humans solve problems using qualitatively different modes of processing.
Kahneman's influential dual-process framework distinguishes a fast, effortless,
heuristic mode (``System~1'') from a slow, effortful, deliberative mode
(``System~2'') \citep{kahneman2011thinking}.  While the binary framing has
faced sustained criticism---Melnikoff and Bargh argue that the ``mythical number
two'' obscures what is better understood as a continuum of processing intensity
\citep{melnikoff2018mythical}---the core observation remains: cognitive systems
flexibly allocate effort, and this allocation has measurable consequences for
accuracy, latency, and susceptibility to systematic biases
\citep{evans2013dual, deneys2012bias}.

Large language models (LLMs) now exhibit strikingly analogous behavioral
patterns.  On classic cognitive bias tasks, instruction-tuned models reproduce
human-like errors on conflict items---problems where a salient heuristic cue
diverges from the normatively correct answer---while succeeding on structurally
matched controls where heuristic and deliberative cues converge
\citep{hagendorff2023human}.  Reasoning-trained models, such as those distilled
from DeepSeek-R1 \citep{deepseekr1}, resist many of these heuristic lures,
producing explicit chain-of-thought (CoT) traces that superficially resemble
deliberative processing.  A growing body of work has documented these behavioral
parallels across models and task families \citep{xie2025cogsci, rao2024cobbler,
codaforno2025dual, ziabari2025reasoning}.

Yet behavior alone cannot distinguish genuine processing-mode differences from
surface mimicry.  A model that memorizes the correct answer to the bat-and-ball
problem produces the same output as one that ``reasons through'' it, but via
radically different internal computation.  Moreover, chain-of-thought traces are
often unfaithful to the underlying computation: recent work shows that the
majority of reasoning steps in CoT---including apparent self-correction
markers---do not causally influence the final answer
\citep{lanham2023measuring, jiang2025aha}.  The critical open question is
therefore not whether LLMs \emph{behave} as if they have fast and slow
processing modes, but whether such modes have \emph{mechanistic correlates} in
the model's internal representations.

Two recent lines of work begin to address this gap.  On the mechanistic side,
probing studies have shown that LLM hidden states encode self-verification
signals before the model generates its answer \citep{zhang2025probing}, that
attention heads specialize for recall versus reasoning
\citep{fartale2025recall}, and that sparse autoencoder (SAE) features can
identify and steer reasoning-related computations
\citep{fang2026sae, yang2026step, goodfire2025r1}.  On the cognitive science
side, CogBias \citep{huang2026cogbias} recently combined linear probes with
activation steering to detect and mitigate cognitive biases in LLM internal
representations.  However, no prior work has attempted a \emph{unified}
mechanistic characterization of dual-process-like computation, combining
multiple complementary interpretability methods on a benchmark explicitly
designed to elicit and control for the distinction.

This paper presents the first such characterization.  We make four
contributions:

\begin{enumerate}[leftmargin=*,itemsep=2pt]
  \item \textbf{A contrastive cognitive bias benchmark} of 284 items (142
  conflict/no-conflict matched pairs) across 7 task categories, with novel
  structural isomorphs designed to resist memorization-based shortcuts
  (Section~\ref{sec:benchmark}).

  \item \textbf{A multi-method mechanistic analysis} combining five
  complementary workstreams---linear probes, SAE feature analysis, attention
  entropy, representational geometry, and causal interventions---applied to four
  models spanning standard and reasoning-trained architectures
  (Section~\ref{sec:methods}).

  \item \textbf{A controlled architecture comparison} between
  Llama-3.1-8B-Instruct and DeepSeek-R1-Distill-Llama-8B, which share
  identical architectures (32 layers, 32 attention heads, 4096 hidden
  dimensions) but differ in training.  Any internal differences are attributable
  to reasoning distillation, not architectural confounds.

  \item \textbf{Built-in falsification} at every stage: Hewitt--Liang control
  tasks for probes \citep{hewitt2019control}, the Ma et al. framework for SAE
  feature validation \citep{ma2026falsification}, random-projection baselines
  for geometry, and random-direction controls for causal interventions.
\end{enumerate}

Critically, we frame our investigation not as testing for a binary System~1
versus System~2 dichotomy, but as measuring the \textbf{deliberation-intensity
gradient}: the degree to which models exhibit mechanistically distinct internal
signatures when processing tasks that differentially recruit heuristic versus
deliberative computation.  This graded framing is both more scientifically
defensible and more informative---it asks not ``does a model have System~2?''
but ``how, where, and when does deliberation intensity vary across the network's
representations?''
